\section{A unified Kimura-family perspective}\label{sec:kimura-perspective}

The completed JC, K2P, and K3P classifications have the same structural
quotient but different local analytic geometry.  The comparison below records
only results proved in the three companion programs
\cite{KriebelJC2026,KriebelK2P2026}; it is not an extrapolation to other
group-based models.

\begin{center}
\small
\begin{tabular}{L{.29\linewidth}ccc}
\toprule
feature & JC & K2P & K3P\\
\midrule
independent nonzero character sectors & 1 & 2 & 3\\
bridge scales per component incidence & 1 & 2 & 3\\
ordinary-triangle normalized local dimension & 4 & 9 & 14\\
triangle image near the common point & ambient & ambient & quartic \(H_{14}\)\\
certified weak-class ambiguity dimension & \(2n+1\), \(n\ge4\) & \(4n-3\), \(n\ge3\) & \(6n-3\), \(n\ge3\)\\
strong-class topology quotient & \(\trianglerel\) & \(\trianglerel\) & \(\trianglerel\)\\
proper directed containment in strong class & none & none & none\\
\bottomrule
\end{tabular}
\end{center}

K3P therefore agrees topologically with the two symmetry-reduced models but
not geometrically.  Its ordinary triangle loses one ambient dimension even
though its three Fourier sectors are independent; the common ambiguity is
the smooth part of an irreducible quartic hypersurface.  Conversely, each
identical cherry exposes six new directions, two per sector, accounting for
the all-\(n\) sharpness dimension.

The common conceptual boundary is strong tree-childness: inside it, positive
bridge factorization, local factor classification, and coherent probes leave
only ordinary triangle redirection; immediately outside it, robust
full-dimensional ambiguities occur.  This synthesis is restricted to JC,
K2P, and K3P on the domains and network class stated in their respective
theorems.  No theorem for arbitrary finite groups, arbitrary state-label
symmetries, or signed Fourier components is asserted.
